\documentclass[a4 paper,fontset=windows]{ctexart}
\usepackage{geometry}
\geometry{top=3cm,bottom=3cm,left=2.5cm,right=2.5cm}

\begin{document}
	\title{\textbf{数据结构与算法K04\ \ \ \ 模拟仿真支持决策}}
	\author{1900011604\ \ \ \ 张植竣\ \ \ \ 指导老师：陈斌\\北京大学物理学院天文学系}
	\date{2020年3月17日}
	\maketitle
	
	\paragraph{定义问题}考虑一个机场，每天有众多航班起飞、降落，但是停机坪和跑道的数量是固定的。怎样优化安排飞机的起降时间和登机口、停机坪、起降跑道的分配，才能使得机场承载的客流量达到最大、跑道拥挤以致于等待起飞降落的时间最短，乘客需要坐摆渡车去远处停机坪登机的概率最小。
	
	\paragraph{各种假设和参数}假设机场有$m$个近机位（用廊桥连接）和$n$个远机位（需要用摆渡车），为简化问题，不考虑停机坪之间除离大楼距离之外的其它性质的差异（如面积、停泊飞机的载客量等）。对于任意一次航班的全部旅客，使用廊桥登机的总时间为$t_1$，使用摆渡车登机的总时间为$t_2\ (t_1<t_2)$。每一架飞机在停机坪检修、加油、机组人员交接班的总时间为$t_0$。机场共有$k$条跑道$(k<m+n)$，且每条跑道只能容纳一架客机起飞或降落，飞机起飞或降落时在跑道上滑行的总时间为$t_{glide}$（不考虑等待时间）。模拟客流量较小、正常、较大三种情况的分配，假设每一架飞机的载客量相同，对应概率分别为$p_1\ flight/h$，$p_2\ flight/h$，$p_3\ flight/h\ (p_1<p_2<p_3)$，这给出了每小时起飞或降落的飞机的数量。	
	
	\paragraph{需要用到的队列}每个跑道都是一个双端队列，即共有$k$个双端队列。飞机起飞定义为从队首退出队列，降落定义为从队首加入队列；飞机离开停机坪进入跑道定义为从队尾加入队列，降落后进入第一个空出的停机坪定义为从队尾退出队列（如果原有停机坪空出则直接退出队列）。按照$p_i\ (i=1,2,3)$的概率在随机时间通知起飞或降落航班，记录每一架飞机在队列内等待的时间$T$和需要在远机位停泊的航班数$N$。
	
	\paragraph{各种概率数据}$m=45,\ n=25,\ k=2,\ t_{glide}=10min,\ t_1=30min,\ t_2=50min,\ p_1=30\ flight/h,\ p_2=40\ flight/h,\ p_3=60\ flight/h$（注：这里的部分数据参考了广州白云机场的数据）
	
\end{document}